\documentclass{article}
\usepackage{/globals/de}
\usepackage{/globals/laws}
\usepackage{/globals/anki}
\ankiprefix{art20}
\ankitopic{Art. 20}
\ankishow
\begin{document}
\section{Artikel 20 \GG} 
\art{20} \GG stellt in vier kurzen Absätzen die Prinzipien des deutschen Staates auf.

\subsection{\absa{1}}
Es wird die deutsche Staatsform festgelegt; eine demokratische Republik. Aus dem Demokratieprinzip geht auch \ua die Rechtsstaatlichkeit hervor.

Deutschland wird des weiteren ein Sozialstaat genannt; die genaue Bedeutung dessen ist aber Interpretationssache und beinhaltet \idr keinen totalen Wohlfartsstaat.

Zudem ist Deutschland föderalistisch; es gibt den Bund, den \sog Gesamtstaat, der aus mehreren Bundesländern, \sog Gliedstaaten, besteht. Der Bund hat generelle Hoheitsrechte, bis auf bei \zb der Polizei- und Bildungspolitik.
\ankinote{abs1}{(1)}{demokratische Republik, Sozialprinzip, Föderalismus}

\subsection{\absa{2}}
Es gibt Wahlen und Abstimmungen. Es gibt die klassische Gewaltenteilung in die Judikative, die Legislative und die Exekutive.
\ankinote{abs2}{(2)}{Gewaltenteilung, Wahlen}

\subsection{\absa{3}}
Verwaltungsakte müssen den Gesetzen folgen, Gesetze müssen der Verfassung folgen.
\ankinote{abs3}{(3)}{Gesetzmäßigkeit; Verwaltung folgt Gesetzen, Gesetze folgen Verfassung}

\subsection{\absa{4}}
Bürger haben ein Widerstandsrecht, gegen die Versuchung die Verfassungsordnung zu beseitigen.
\ankinote{abs4}{(4)}{Widerstandsrecht}

\end{document}